\begin{recipe}{Salade met honing-mosterddressing}
\kicker[Bijgerecht,Vegetarisch]{Bijgerecht · vegetarisch}
\index[register]{Salade met honing-mosterddressing}
\index[register]{Rozijnen}
\index[register]{Komkommer}
\index[register]{Paprika}
\index[register]{Wortel}

\meta{15 minuten}

\lettrine{E}{en} frisse salade met sla, rozijnen, komkommer, paprika en
wortel, afgemaakt met een romige dressing van maar vijf ingrediënten. De
zoete rozijnen en de honing-mosterddressing maken dit ook voor kinderen een
laagdrempelige manier om groente te eten.

\blockrule{Ingrediënten}
\begin{ingredients}
  \ing{150 g}{gemengde sla}
  \ing{50 g}{rozijnen}\index[register]{Rozijnen}
  \ing{½}{komkommer}\index[register]{Komkommer}
  \ing{½}{paprika}\index[register]{Paprika}
  \ing{1 kleine}{wortel}\index[register]{Wortel}
  \ing{3 el}{olijfolie}
  \ing{1 tl}{honing}
  \ing{1 tl}{mosterd}
  \ingb{peper en zout, een snufje van elk}
\end{ingredients}

\blockrule{Bereiding}
\begin{steps}
  \item Was de komkommer, de paprika en de wortel. Snijd de komkommer en de
        paprika in blokjes en rasp de wortel, of snijd hem in dunne
        reepjes.
  \item Week de rozijnen 10 minuten in een kommetje water en giet ze goed
        af.
  \item Klop in een kommetje de olijfolie, honing, mosterd, peper en zout
        krachtig door elkaar tot een gladde, romige dressing.
  \item Doe de sla, rozijnen, komkommer, paprika en wortel in een grote
        schaal. Giet de dressing erover en hussel het geheel goed door
        elkaar.
\end{steps}
\end{recipe}
